% Generated draft: requires mathematical and source review.
\documentclass[11pt]{article}
\usepackage[margin=25mm]{geometry}
\usepackage{fontspec}
\setmainfont{Latin Modern Roman}
\usepackage{amsmath,amssymb,mathtools,unicode-math}
\setmathfont{Latin Modern Math}
\usepackage{enumitem,needspace,xurl,hyperref}
\setlength{\parindent}{0pt}
\setlength{\parskip}{5pt}
\setlength{\emergencystretch}{3em}
\newcommand{\N}{\mathbb{N}}
\newcommand{\Z}{\mathbb{Z}}
\newcommand{\Q}{\mathbb{Q}}
\newcommand{\R}{\mathbb{R}}
\newcommand{\C}{\mathbb{C}}
\newcommand{\F}{\mathbb{F}}
\newcommand{\A}{\mathbb{A}}
\newcommand{\PP}{\mathbb{P}}
\newcommand{\E}{\mathbb{E}}
\newcommand{\eps}{\varepsilon}
\newcommand{\dd}{\,\mathrm{d}}
\begin{document}
\section*{UnsolvedMath ID: 10800010}
\textbf{Source problem code:} catalogue entry
\subsection{Definitions and mathematical statement}
Problem 4 asks: ``Are these geometric obstructions sufficient to solve the above problem?'' Both the set of geometric obstructions and the preceding problem are missing from the supplied statement. No precise sufficiency implication can therefore be specified without recovering those antecedents.

The supplied research summary places the question in Arnold's algebraic-volume problem and interprets the literal sufficiency issue as existence of an obstructing singularity type on every relevant complexification. It does not define the domains, ambient dimension, complexification, singularity types or obstruction theorem. For terminology only, a common volume-function setting assigns to a domain \(D\subset\R^n\), a linear functional \(\ell\), and a cut level \(t\), the volume
\[
 V_D(\ell,t)=\operatorname{vol}_n\bigl(D\cap\{x:\ell(x)\le t\}\bigr).
\]
This explains the topic without selecting the omitted antecedent or an algebraicity convention for its analytic continuation.

The summary reports a restricted ellipsoid classification attributed to Agranovsky (2017), for smooth boundaries with algebraic volume function under an essential no-real-singular-points hypothesis. It also reports Vassiliev (2020) candidates with finitely valued analytic continuation and says the full classification remains open. It explicitly says the literal geometric-obstruction sufficiency question was not directly addressed in the literature it could verify. Those qualifications are retained: the restricted classification is not a solution of the undefined general implication, and candidates are not asserted to be verified counterexamples.
\subsection{Short English statement}
Determine whether the source's previously described geometric obstructions settle its preceding algebraic-volume problem. The antecedents are absent. The supplied partial-progress discussion distinguishes a restricted ellipsoid result from the literal sufficiency question and an unresolved full classification.
\subsection{Sources}
\begin{itemize}
\item Local UnsolvedMath registry, permanent ID \(10800010\), exact source code \texttt{catalogue entry}, Problem 4, attributed to Victor Vassiliev; no proposed year supplied.
\item Vassiliev, \emph{A Few Problems on Monodromy and Discriminants}, Problem 4: \url{https://amj.math.stonybrook.edu/html-articles/Files-2015-2024/15-11/}. The background describes eighteen explicitly labeled theorem-style problem environments.
\item Supplied research summary names Agranovsky (2017) for a restricted ellipsoid classification and Vassiliev (2020) for finitely valued analytic-continuation candidates; no titles, identifiers or complete references are supplied. It expressly reports no direct verification of the literal sufficiency question.
\item HTML-structure extraction, accessed 2026-07-29. Source evidence NEEDS\_REVIEW presents an open problem and requires a current primary-literature audit. Registry status partially\_solved and classification PARTIAL-PROGRESS retain the restricted, qualified assessment above.
\item Dataset attribution: Ulam AI / UnsolvedMath Contributors, CC BY 4.0. Source rights NEEDS\_REVIEW: publicly accessible author/publisher copy, redistribution terms requiring review.
\end{itemize}
\end{document}
